%%%%%%%%%%%%%%%%%%%%%%%%%%%%%%%%%%%%%%%%%%%%%%%%%%%%%%%%%%%
% IEEE Transactions / Journal Template
% --------------------------------------------------------
% 出版社: IEEE (Institute of Electrical and Electronics Engineers)
% 适用对象: IEEE Transactions, IEEE Journals, IEEE Letters
% 编译方式: pdflatex Main_IEEE_Trans -> bibtex Main_IEEE_Trans -> pdflatex Main_IEEE_Trans -> pdflatex Main_IEEE_Trans
%%%%%%%%%%%%%%%%%%%%%%%%%%%%%%%%%%%%%%%%%%%%%%%%%%%%%%%%%%%

\documentclass[journal,10pt,twocolumn]{../../packages/publishers/ieee/IEEEtran}

\usepackage{../../packages/publishers/compat_helper}
\usepackage{cite}
\usepackage{algorithmic}
\usepackage{textcomp}
\graphicspath{{../../figures/}}

% 墨绿色学术高亮配置
\usepackage{hyperref}
\hypersetup{
    colorlinks = true,
    linkcolor  = academicgreen,
    citecolor  = academicgreen,
    urlcolor   = academicgreen
}

% ---------------------------------------------------------
% 论文元数据 (Metadata)
% ---------------------------------------------------------
\title{Modular Academic LaTeX Framework: Seamless Multi-Publisher Manuscript Migration}

\author{Hongquan~Zhou,~\IEEEmembership{Graduate~Student~Member,~IEEE,}
        Xiaolin~Jia,~\IEEEmembership{Senior~Member,~IEEE,}
        and~Zhong~Du,~\IEEEmembership{Member,~IEEE}%
\thanks{This work was supported in part by the National Natural Science Foundation of China.}%
\thanks{H. Zhou and X. Jia are with the Mobile Internet of Things and RFID Technology Key Laboratory, Southwest University of Science and Technology, Mianyang 621010, China (e-mail: my\_jiaxl@163.com).}%
\thanks{Z. Du is with the School of Computer Science and Technology, Southwest University of Science and Technology, Mianyang 621010, China.}%
}

\markboth{IEEE Transactions on Knowledge and Data Engineering,~Vol.~XX, No.~X, August~2026}%
{Zhou \MakeLowercase{\textit{et al.}}: Modular Academic LaTeX Framework}

\begin{document}

\maketitle

% ---------------------------------------------------------
% 摘要与关键词
% ---------------------------------------------------------
\begin{abstract}
Tags identification in dense mobile RFID systems presents a major challenge due to the stochastic residence time of tags and severe communication collisions. This paper demonstrates a unified decoupled academic authoring framework that bridges modern modular writing with strict IEEE Transactions publishing requirements. Core sections are maintained independently from venue-specific formatting, ensuring instantaneous multi-target manuscript compilation and verification.
\end{abstract}

\begin{IEEEkeywords}
RFID systems, anti-collision protocol, IEEE Transactions, modular authoring, reproducible research.
\end{IEEEkeywords}

% ---------------------------------------------------------
% 共享正文内容模块
% ---------------------------------------------------------
% ==========================================================
% Section: Introduction
% ==========================================================

\section{Introduction}
\label{sec:introduction}

\dropcapstart{W}ith the rapid advancement of modern Internet of Things (IoT) ecosystems and cyber-physical infrastructure, automated Radio Frequency Identification (RFID) systems have gained immense traction across supply chains, smart manufacturing, and logistics~\cite{Industry5, RFID-application1}. Fast and reliable identification of dense dynamic tags is essential for maintaining high throughput and mission-critical tracking fidelity.

\subsection{Research Background and Motivation}
In industrial environments, tags move dynamically across interrogation zones under stochastic residence intervals~\cite{RFID-in-IOT}. Conventional anti-collision protocols based on static frame-slotted ALOHA or deterministic tree splitting encounter severe communication latency and tag starvation under heavy workload surges~\cite{EPC-C1-Gen2, TAD}. Consequently, architecting a unified, modular framework that decouples the scientific algorithmic representation from venue-specific typography is of substantial benefit to reproducible research and rapid dissemination.

\subsection{Primary Contributions}
The major contributions of this work are three-fold:
\begin{itemize}
    \item We design a decoupled multi-publisher template architecture that isolates manuscript content from target venue typographical dependencies.
    \item We seamlessly integrate official journal standards (including IEEE Transactions, Elsevier CAS, and ACM SIGCONF) with full support for authentic bibliographic databases.
    \item We empirically validate seamless template switching without modifying the underlying technical discourse or experimental data.
\end{itemize}

% ==========================================================
% Section: Related Work
% ==========================================================

\section{Related Work}
\label{sec:related_work}

Extensive investigations have been conducted on anti-collision protocols and academic authoring pipelines.

\subsection{RFID Anti-Collision Protocols}
Tag identification protocols are traditionally categorized into ALOHA-based probabilistic schemes and tree-based deterministic algorithms~\cite{EPC-C1-Gen2}. While probabilistic methods dynamically adjust frame length based on tag cardinality estimation, tree-splitting methods resolve collisions by recursively partitioning the tag population into collision-free subsets~\cite{TAD}. Recent developments incorporate physical-layer signal decoding and predictive scheduling to alleviate collision overhead~\cite{RFID-in-IOT}.

\subsection{Modular Document Synthesis}
In scientific publication, managing submissions across divergent publishing ecosystems—such as IEEE, Elsevier, and ACM—often introduces substantial friction due to incompatible author markup, differing citation models (e.g., BibTeX vs. BibLaTeX), and layout constraints. Developing a modular, decoupled framework eliminates redundant adaptations and enhances research reproducibility.

% ==========================================================
% Section: Methodology
% ==========================================================

\section{Methodology}
\label{sec:methodology}

In this section, we formulate the decoupled template framework and its mathematical underpinning for content reusability.

\subsection{Mathematical Formulation}
Let $\mathcal{C}$ represent the set of core scientific content sections, and let $\mathcal{T} = \{T_{\text{custom}}, T_{\text{ieee}}, T_{\text{elsevier}}, T_{\text{acm}}, T_{\text{springer}}\}$ denote the set of target publisher templates. A compiled manuscript $\mathcal{M}_i$ for a given publisher template $T_i \in \mathcal{T}$ is obtained through a deterministic mapping $\Phi$:
\begin{equation}
    \mathcal{M}_i = \Phi(T_i, \mathcal{C}, \mathcal{B}, \mathcal{F})
    \label{eq:manuscript_mapping}
\end{equation}
where $\mathcal{B}$ denotes the bibliography database and $\mathcal{F}$ represents the visual figure assets.

\subsection{Architecture Overview}
The system enforces strict invariance on $\mathcal{C}$:
\begin{equation}
    \frac{\partial \mathcal{C}}{\partial T_i} = 0, \quad \forall T_i \in \mathcal{T}
\end{equation}
This invariance ensures that modifying typography or publisher requirements never introduces discrepancies or side effects into the scientific discourse.

% ==========================================================
% Section: Experiments and Results
% ==========================================================

\section{Experiments and Results}
\label{sec:experiments}

To evaluate the efficiency of the multi-publisher template suite, we benchmark migration overhead, compilation time, and typographic compliance across diverse venue standards.

\subsection{Experimental Setup}
All tests are conducted under a unified TeX Live distribution. We compare the migration effort measured in required manual modifications against traditional unified or duplicated repository patterns.

\begin{table}[htbp]
    \centering
    \caption{Comparison of Manuscript Management Strategies}
    \label{tab:comparison}
    \begin{tabularx}{\linewidth}{Y Y Y Y}
        \toprule
        \textbf{Strategy} & \textbf{Separation Level} & \textbf{Switching Cost} & \textbf{Conflict Risk} \\
        \midrule
        Monolithic Multi-Repo & Low & High & High \\
        Conditional Engine & Medium & Medium & High \\
        \textbf{Decoupled (Ours)} & \textbf{High} & \textbf{Minimal} & \textbf{None} \\
        \bottomrule
    \end{tabularx}
\end{table}

\subsection{Performance and Compatibility}
As summarized in \tabref{tab:comparison}, the proposed architecture eliminates package collision hazards while keeping manuscript synchronization effortless.

% ==========================================================
% Section: Conclusion
% ==========================================================

\section{Conclusion}
\label{sec:conclusion}

In this paper, we introduced a modular, multi-publisher LaTeX template ecosystem designed for streamlined academic authoring and painless journal/conference migration. By decoupling core technical content from venue-specific document styles, researchers can seamlessly submit, reformat, and archive their manuscripts with minimal friction. Future work includes expanding automated verification scripts and continuous integration pipelines for collaborative paper authoring.


% ---------------------------------------------------------
% 参考文献 (前置 journal_abrv 自动缩写)
% ---------------------------------------------------------
\bibliographystyle{../../packages/publishers/ieee/IEEEtran}
\bibliography{../../bib/journal_abrv,../../bib/refer}

\end{document}
